% =====================================================================
\section{Mathematical foundations, and the validation they had to pass}
\label{sec:measures}
% =====================================================================

This section is the reference core of the briefing: every quantity the study
computes, as a formula, with the source it originates from and the
implementation detail that makes it correct on this exchange. The full
variable list is in Appendix~\ref{app:variables}; the ladder algorithm has its
own appendix. Throughout, $i$ indexes stocks, $t$ trading days, and prices are
handled internally as integers in tenths of a yen, so the modular arithmetic
below is exact on every grid --- floating-point \texttt{mod} silently corrupts
the 0.1-yen grid, which is exactly where clustering is strongest.

\subsection{Clustering (Ohta 2026)}

For tick size $\tau$, the \emph{last price digit} of a trade at price $p$ is
%
\begin{equation}
d(p;\tau) \;=\; \frac{p \bmod 10\tau}{\tau} \;\in\; \{0,1,\dots,9\},
\label{eq:digit}
\end{equation}
%
and a \emph{round price} is one with $d(p;\tau)=0$
\parencite{Ohta2026NoiseTraders}. The construction requires the tick to be a
power of ten --- the paper's sample filter (d) --- and under a uniform digit
distribution each digit carries 10\% of volume, which is the benchmark every
clustering number in this study is read against. The stock-day measures are
volume-weighted shares at the round price,
%
\begin{equation}
M^{0}_{i,t} \;=\;
\frac{\sum_{\text{trades}} v \,\mathbb{1}[d(p;\tau)=0]}{\sum_{\text{trades}} v},
\qquad
M^{k\,\mathit{Size}\,0}_{i,t} \;=\;
\frac{\sum_{k,\mathit{Size}} v \,\mathbb{1}[d(p;\tau)=0]}
     {\sum_{k,\mathit{Size}} v},
\label{eq:M}
\end{equation}
%
with initiator $k \in \{B, S\}$ from the exchange's own execution-type field
(so no \textcite{LeeReady1991} inference is needed), \emph{Large} meaning more
than one trading unit and \emph{Small} exactly one, following the paper. Only
continuous-session (zaraba) trades executed against a two-sided ordinary quote
enter; auction prints are excluded because their prices are not set by
individual limit orders. The determinants of clustering --- volatility, price
level, capitalisation, trading frequency, tick --- come from
\textcite{Harris1991Clustering}, which is where the control set originates.

\subsection{Liquidity (the paper's definitions)}

With $m$ the midquote prevailing immediately before a trade and $q=+1$
($-1$) for buy-initiated (sell-initiated) trades, the \emph{effective
half-spread}, \emph{price impact} at horizon $\Delta$, and \emph{realized
spread} are the volume-weighted means of
%
\begin{equation}
\mathit{es} = q\,\frac{p - m_t}{m_t}, \qquad
\mathit{imp}_\Delta = q\,\frac{m_{t+\Delta} - m_t}{m_t}, \qquad
\mathit{rs}_\Delta = \mathit{es} - \mathit{imp}_\Delta,
\label{eq:spread}
\end{equation}
%
in basis points, at $\Delta \in \{1\mathrm{s}, 60\mathrm{s}, 300\mathrm{s}\}$.
The decomposition of the effective spread into impact plus realized spread
follows \textcite{HendershottJonesMenkveld2011}; reading impact as the
informational content of a trade goes back to \textcite{GlostenMilgrom1985}
and \textcite{Hasbrouck1991}, and Ohta's point --- adopted here --- is that in
a limit-order market impact also rises when stale orders sit waiting to be
taken, which is not information. Note there is no factor of two: these are
half-spreads, matching the paper's convention. A trade's horizon must end
inside its own session (no impact is measured across the lunch break or the
close), and $\mathit{RS} = \mathit{ES} - \mathit{Imp}$ holds as an identity on
each horizon's own trade set --- verified to three decimals as an
implementation check, not a finding. The round-price impact premium is
%
\begin{equation}
\Delta\mathit{Imp}^{k\mathit{Large}}_{i,t} \;=\;
\mathit{Imp}^{k\mathit{Large},\,d=0}_{i,t} -
\mathit{Imp}^{k\mathit{Large},\,d\neq 0}_{i,t},
\label{eq:dimp}
\end{equation}
%
with cells of fewer than five trades set missing rather than reported.
\textcite{BoxGriffith2016} is the closest antecedent for pairing clustering
asymmetries with round-price impact. The auxiliary Amihud measure is
$|r^{oc}_{i,t}|$ over turnover in millions of yen \parencite{Amihud2002}.

\subsection{Price formation}

Realised variance is the sum of squared five-minute midquote log returns
within each session, and the variance-ratio deviation is
%
\begin{equation}
\left| \mathit{VR}_{i,t} - 1 \right|, \qquad
\mathit{VR}_{i,t} = \frac{\sigma^2_{5\mathrm{min}}}{5\,\sigma^2_{1\mathrm{min}}},
\label{eq:vr}
\end{equation}
%
the standard variance-ratio construction measuring departure from a random
walk in either direction; returns never span the lunch break or the overnight
gap.

\subsection{The visible ten-level book}

Two measures use data the paper did not have. The \emph{round-price resting
depth} is the share of visible ten-level depth standing at round prices,
%
\begin{equation}
\mathit{RDepth}^{\mathrm{ask}0}_{i,t} \;=\;
\text{time-avg}\;
\frac{\sum_{\ell \le 10} v^{\mathrm{ask}}_{\ell}\,
      \mathbb{1}[d(p^{\mathrm{ask}}_{\ell};\tau) = 0]}
     {\sum_{\ell \le 10} v^{\mathrm{ask}}_{\ell}},
\label{eq:rdepth}
\end{equation}
%
sampled on a one-minute grid over ordinary two-sided snapshots (uniform
benchmark again 10\%). The idea that the book itself is thicker at round
prices is documented for Hong Kong by \textcite{AhnCaiCheung2005}; the measure
here is that observation turned into a stock-day share. The \emph{order-flow
imbalance} at the best quote follows \textcite{ContKukanovStoikov2014}
exactly, and the \emph{depth slope} --- how fast size grows walking away from
the best quote --- follows \textcite{NaesSkjeltorp2006};
\textcite{CaoHanschWang2009} is the license for looking past the best quote at
all.

\subsection{Limit-order flow inferred from the ladder}

Ohta's explanatory variables are the round-price shares of submitted and
cancelled limit-order volume. The feed does not carry orders, so they are
inferred from the accounting identity that resting volume at a price can only
change by execution, cancellation or submission: with
$\delta = V_n(p) - V_{n-1}(p)$ between consecutive snapshots and $X_n(p)$ the
volume printed at $p$ in the interval,
%
\begin{equation}
\mathit{submitted}_n(p) = \max\{\delta + X_n(p),\, 0\}, \qquad
\mathit{cancelled}_n(p) = \max\{-(\delta + X_n(p)),\, 0\}.
\label{eq:ladder}
\end{equation}
%
$L^{k0}$ (round-price share of submissions) and $L^{k0C}$ (cancellations over
submissions at round prices) aggregate these to the stock-day, targeting the
same quantities the paper builds from its feed. Two rules make
equation~\eqref{eq:ladder} correct in practice --- match by price, never by
level index, and count a price only where it is observable in both snapshots
--- and Appendix~\ref{app:ladder} states them, the three known failure modes,
and the hand-built order book that tests all three.

\subsection{Estimation and inference}

The daily specifications are two-way fixed-effects panels,
%
\begin{equation}
y_{i,t} \;=\; \beta\, M^{B\mathit{Large}0}_{i,t-1}
\;+\; \gamma\, y_{i,t-1}
\;+\; \boldsymbol{\theta}' \mathbf{x}_{i,t-1}
\;+\; \alpha_i \;+\; \delta_t \;+\; \varepsilon_{i,t},
\label{eq:fe}
\end{equation}
%
with standard errors clustered by stock and by day following
\textcite{CameronGelbachMiller2011} --- with the same stock appearing on
nearly every day and the same day covering a thousand stocks, one-way
clustering is not conservative, it is wrong. The control vector
$\mathbf{x}$ holds the relative tick, lagged log turnover, lagged log realised
variance, lagged log spread, the overnight return, and the opening-digit
contamination interactions of Section~\ref{sec:stylized}, timed to match the
regressor. Including the lagged outcome makes \eqref{eq:fe} a dynamic panel;
the induced bias is of order $1/T$ \parencite{Nickell1981} and negligible at
$T$ near eighty, and the robustness section re-estimates with five lags so
that ``the outcome's own history'' does not secretly mean one day. The
estimator cross-check runs the same regression cross-sectionally day by day
and averages the coefficients \parencite{FamaMacBeth1973} with
\textcite{NeweyWest1987} standard errors over the daily series; agreement or
disagreement between \eqref{eq:fe} and that average is informative about
whether a relationship lives within stocks over time or between stocks, and
Section~\ref{sec:robust} reports which it is here. Continuous variables are
winsorised within stock at the 2.5th and 97.5th percentiles, the paper's
convention; the intraday specifications replace $\alpha_i + \delta_t$ with
stock-day and bucket-of-day effects and cluster by stock and by date.

\subsection{What the measurement layer had to prove}

The measure library is a single module, and the full panel build calls it
unchanged; validating it is validating production code. It had to pass four
things.

\paragraph{Two hand-verified stock-days, to the second decimal.}
Two stock-days were computed independently in advance and used as tripwires: one
on the 1-yen grid, one on the 0.1-yen grid where the digit arithmetic is fragile.
All ten clustering cells reproduce to within 0.005 percentage points, and both
tick sizes were recovered identically by the coded exchange schedule and by
inference from the tape.

\paragraph{A negative control.}
Excluding the opening and closing auctions is not cosmetic. Letting them back in
moves the daily measure by $-2.91$ and $+8.73$ percentage points on the two
anchor days: the closing call is a single enormous volume-weighted print, and its
price is not set by an individual limit order the way continuous-session prices
are. A pipeline that silently included it would produce measures that looked
plausible and were wrong.

\paragraph{A hand-built order book.}
The ladder inference is checked against a six-snapshot tape whose order flow is
known exactly, constructed so that each of the three standard failure modes
produces a visibly wrong answer --- mistaking an execution for a cancellation,
inventing flow when the ladder shifts a level, and counting depth that entered
from beyond the visible window. Appendix~\ref{app:ladder} gives the construction.

\paragraph{Magnitudes nobody tuned for --- reported with their one miss.}
The strongest check is external. Ohta reports a mean round-price share of
submitted sell limit orders of 0.106 and a mean cancellation ratio of 0.806
over 2010--2022. Run on this study's panel, the inference lands on the
submission share to half a percentage point without calibration --- and
overshoots the cancellation ratio (Appendix~\ref{app:ladder} carries the
numbers). The agreement on levels is treated as validation; the overshoot on
the ratio is treated as an open difference, not absorbed into the claim.

\paragraph{Verdict.}
The measurement layer reproduces two independently computed stock-days exactly,
satisfies the identities it should satisfy, orders itself correctly across the
liquidity spectrum, and lands on the published submission-share magnitudes it
was never shown. The results that follow rest on it.
